\chapter{Probability and Gaussian Algebra}
\label{ch:prob}

\begin{objectives}
\begin{itemize}
  \item Propagate means and covariances through linear and linearized maps.
  \item Marginalize and condition Gaussians in both covariance and information form, and know which
        operation is cheap in which form.
  \item Connect Bayesian estimation, maximum a posteriori estimation and weighted least squares.
  \item Understand the $\chi^2$ distribution that underlies gating and consistency tests.
  \item Model continuous-time sensor noise and random walks, and derive how integrated noise grows.
\end{itemize}
\end{objectives}

A filter is a machine for maintaining a probability distribution over the state. Kalman filters
restrict that distribution to be Gaussian, which reduces the whole machinery to the manipulation of
means and covariance matrices. This chapter develops that algebra.

% ------------------------------------------------------------------------------
\section{Random vectors, mean and covariance}

For a random vector $x \in \R^n$ with density $p(x)$, the mean and covariance are
\begin{equation}
  \mu = \E[x] = \int x\, p(x)\, \dd x, \qquad
  \Sigma = \Cov(x) = \E\big[(x - \mu)(x - \mu)\T\big].
\end{equation}
$\Sigma$ is symmetric and positive semi-definite: for any $a$,
$a\T\Sigma a = \E[(a\T(x - \mu))^2] \ge 0$. Its diagonal holds the variances $\sigma_i^2$ and its
off-diagonal entries the covariances. The correlation coefficient
$\rho_{ij} = \Sigma_{ij}/(\sigma_i\sigma_j) \in [-1, 1]$ measures how strongly two components move
together. The cross-covariance of two random vectors is
$\Sigma_{xy} = \E[(x - \mu_x)(y - \mu_y)\T]$.

\begin{proposition}[Affine transformations]
\label{prop:affine}
If $y = Ax + b$ then $\E[y] = A\mu + b$ and $\Cov(y) = A\Sigma A\T$. If $x$ and $w$ are uncorrelated,
$\Cov(Ax + Bw) = A\Sigma_x A\T + B\Sigma_w B\T$.
\end{proposition}
\begin{proof}
Linearity of expectation gives the mean. Then
$\Cov(y) = \E[A(x - \mu)(x - \mu)\T A\T] = A\Sigma A\T$. With uncorrelated $x$ and $w$, the cross terms
vanish.
\end{proof}

This one proposition is the covariance propagation of the Kalman filter
($P \leftarrow \Phi P \Phi\T + Q$), the covariance of a stacked measurement, and the augmentation
formula for stochastic cloning ($P \leftarrow JPJ\T$).

% ------------------------------------------------------------------------------
\section{The multivariate Gaussian}

\begin{definition}[Gaussian distribution]
$x \sim \N(\mu, \Sigma)$ with $\Sigma \succ 0$ has density
\begin{equation}
  p(x) = \frac{1}{\sqrt{(2\pi)^n \det\Sigma}}
         \exp\!\Big( -\tfrac12 (x - \mu)\T \Sigma\inv (x - \mu) \Big).
\end{equation}
The quantity $d_M^2(x) = (x - \mu)\T\Sigma\inv(x-\mu)$ is the squared \emph{Mahalanobis distance}
of $x$ from the mean.
\end{definition}

The Mahalanobis distance measures distance in units of standard deviations along each principal
axis, so it is the natural way to ask whether a vector is ``far'' from a distribution. The level
sets $d_M^2 = c$ are the ellipsoids of the previous chapter.

The Gaussian is closed under every operation a Kalman filter performs. Affine maps of Gaussians are
Gaussian, and marginals, conditionals and products of Gaussian densities are Gaussian. This is
why a Gaussian belief can be propagated exactly through a linear system with Gaussian noise, and
why the Kalman filter is exact in that case.

\paragraph{Sampling.} To draw from $\N(\mu,\Sigma)$, factor $\Sigma = LL\T$ (Cholesky), draw
$s \sim \N(0, I)$, and return $\mu + L s$. By \cref{prop:affine} its covariance is $LL\T = \Sigma$.
Monte Carlo tests of covariance propagation (Milestone 5) use exactly this.

\paragraph{Information form.} A Gaussian can equally be parameterized by the \emph{information
matrix} $\Lambda = \Sigma\inv$ and the \emph{information vector} $\eta = \Sigma\inv\mu$. Then
$-\log p(x) = \tfrac12 x\T\Lambda x - \eta\T x + \text{const}$. Information form is natural for
combining independent pieces of evidence, since their information adds, and it is the form that
optimization-based estimators work in.

% ------------------------------------------------------------------------------
\section{Marginals and conditionals}

Partition a jointly Gaussian vector as
\begin{equation}
  \begin{bmatrix} a \\ b \end{bmatrix} \sim
  \N\!\left( \begin{bmatrix} \mu_a \\ \mu_b \end{bmatrix},
  \begin{bmatrix} \Sigma_{aa} & \Sigma_{ab} \\ \Sigma_{ba} & \Sigma_{bb} \end{bmatrix} \right),
  \qquad
  \Lambda = \Sigma\inv = \begin{bmatrix} \Lambda_{aa} & \Lambda_{ab} \\ \Lambda_{ba} & \Lambda_{bb} \end{bmatrix}.
\end{equation}

\begin{theorem}[Gaussian marginalization and conditioning]
\label{thm:gauss-cond}
\begin{enumerate}
  \item The marginal of $a$ is $\N(\mu_a, \Sigma_{aa})$.
  \item The conditional of $a$ given $b$ is $\N(\mu_{a|b}, \Sigma_{a|b})$ with
  \begin{equation}
    \mu_{a|b} = \mu_a + \Sigma_{ab}\Sigma_{bb}\inv (b - \mu_b), \qquad
    \Sigma_{a|b} = \Sigma_{aa} - \Sigma_{ab}\Sigma_{bb}\inv\Sigma_{ba}.
    \label{eq:gauss-cond}
  \end{equation}
\end{enumerate}
\end{theorem}
\begin{proof}
Define $\tilde a = a - \Sigma_{ab}\Sigma_{bb}\inv b$. Then
$\Cov(\tilde a, b) = \Sigma_{ab} - \Sigma_{ab}\Sigma_{bb}\inv\Sigma_{bb} = 0$, so $\tilde a$ and $b$ are
uncorrelated, and being jointly Gaussian they are independent. By \cref{prop:affine},
$\tilde a \sim \N(\mu_a - \Sigma_{ab}\Sigma_{bb}\inv\mu_b,\ \Sigma_{aa} - \Sigma_{ab}\Sigma_{bb}\inv\Sigma_{ba})$.
Conditioning on $b$ leaves the distribution of $\tilde a$ unchanged, and $a = \tilde a + \Sigma_{ab}\Sigma_{bb}\inv b$,
which gives \eqref{eq:gauss-cond}. The marginal follows from \cref{prop:affine} with the selection
matrix $A = [I\ \ 0]$.
\end{proof}

The conditional covariance $\Sigma_{a|b}$ is the Schur complement $\Sigma/\Sigma_{bb}$. Applying
\cref{prop:block-inverse} to $\Sigma$ shows that $\Lambda_{aa} = \Sigma_{a|b}\inv$ and that the marginal
information matrix is a Schur complement of $\Lambda$. The two parameterizations therefore swap the
roles of the two operations:

\begin{center}\small
\begin{tabular}{@{}lll@{}}
\toprule
 & covariance form $(\mu, \Sigma)$ & information form $(\eta, \Lambda)$ \\
\midrule
marginalize out $b$ & $\Sigma_{aa}$ \hfill(\textbf{delete rows/columns}) & $\Lambda_{aa} - \Lambda_{ab}\Lambda_{bb}\inv\Lambda_{ba}$ \hfill(Schur compl.) \\
condition on $b$ & $\Sigma_{aa} - \Sigma_{ab}\Sigma_{bb}\inv\Sigma_{ba}$ \hfill(Schur compl.) & $\Lambda_{aa}$ \hfill(\textbf{delete rows/columns}) \\
\bottomrule
\end{tabular}
\end{center}

\begin{keyidea}[marginalization is free in covariance form]
A Kalman filter stores covariances, so discarding a variable (the oldest clone of the MSCKF window)
is literally deleting its rows and columns from $P$. Optimization-based estimators store
information matrices, so the same operation requires a Schur complement, which creates fill-in and
the ``marginalization prior'' of sliding-window smoothers.
\end{keyidea}

% ------------------------------------------------------------------------------
\section{Bayesian estimation}

Estimation combines a \emph{prior} $p(x)$ with a \emph{likelihood} $p(z\mid x)$, the probability
of the measurement given the state, through Bayes' rule:
\begin{equation}
  p(x \mid z) = \frac{p(z\mid x)\, p(x)}{p(z)} \propto p(z \mid x)\, p(x).
\end{equation}
The posterior $p(x\mid z)$ is the complete answer. Point estimates summarize it:
\begin{itemize}
  \item the \emph{minimum mean-square error} (MMSE) estimate $\hat x = \E[x \mid z]$ minimizes
        $\E[\|x - \hat x\|^2 \mid z]$;
  \item the \emph{maximum a posteriori} (MAP) estimate maximizes $p(x \mid z)$;
  \item the \emph{maximum likelihood} (ML) estimate maximizes $p(z\mid x)$ and ignores the prior.
\end{itemize}

\begin{proposition}[Linear-Gaussian case]
\label{prop:map-ls}
If $x \sim \N(\mu, P)$ and $z = Hx + n$ with $n \sim \N(0, R)$ independent of $x$, the posterior is
Gaussian, MMSE and MAP coincide, and the estimate solves the weighted least-squares problem
\begin{equation}
  \hat x = \argmin_x\ \|x - \mu\|^2_{P\inv} + \|z - Hx\|^2_{R\inv},
  \qquad \|e\|^2_{W} := e\T W e.
  \label{eq:map-ls}
\end{equation}
\end{proposition}
\begin{proof}
$-\log p(x\mid z) = \tfrac12\|x - \mu\|^2_{P\inv} + \tfrac12\|z - Hx\|^2_{R\inv} + \text{const}$
is a convex quadratic in $x$, so the posterior is Gaussian. A Gaussian's mode equals its mean.
\end{proof}

\Cref{ch:kf} solves \eqref{eq:map-ls} in closed form and obtains the Kalman update. The same
objective with nonlinear $h(x)$ in place of $Hx$, iterated over a window of states, is the cost
function of optimization-based VIO.

% ------------------------------------------------------------------------------
\section{Nonlinear transformations and linearization}
\label{sec:linearization}

If $y = f(x)$ with $f$ nonlinear and $x \sim \N(\mu, \Sigma)$, then $y$ is generally \emph{not}
Gaussian. The extended Kalman filter approximates it with a Gaussian by linearizing $f$ about the mean:
\begin{equation}
  y \approx f(\mu) + J(x - \mu), \quad J = \frac{\partial f}{\partial x}\Big|_{\mu}
  \qquad\Longrightarrow\qquad
  y \mathrel{\dot\sim} \N\big(f(\mu),\ J\Sigma J\T\big).
\end{equation}
The approximation is good when $f$ is nearly linear over the region where $x$ has most of its
probability mass, a few standard deviations around $\mu$. A second-order expansion shows that
curvature also biases the mean by approximately $\tfrac12\sum_i e_i \tr(\nabla^2 f_i\, \Sigma)$, which
first-order linearization ignores.

\begin{figure}[t]
  \centering
  \includegraphics[width=0.84\textwidth]{fig_banana.pdf}
  \caption{A robot drives 10~m with a heading uncertainty of $\sigma_\psi = 0.25$~rad. The true
  distribution of its final position (samples) is a curved ``banana''. The linearized Gaussian
  (ellipses) matches its first two moments only locally, puts probability where the robot cannot be,
  and misses the curvature. Errors of this kind are why an estimator must keep its errors small, the
  central design principle of the error-state filter.}
  \label{fig:banana}
\end{figure}

\Cref{fig:banana} shows a classic failure. When heading uncertainty is large, positions are
distributed along an arc, and the linearized Gaussian is a poor description. Two remedies are
available. One is to reduce the uncertainty that must be linearized: the error-state filter of
\cref{ch:eskf} keeps the estimation error, not the state, as the quantity it linearizes, and
continually resets it to zero. The other is to choose coordinates in which the distribution
\emph{is} close to Gaussian. A Gaussian on the heading and travelled distance, pushed through the
exact kinematics, reproduces the banana. \Cref{ch:lie} develops that idea for rotations.

% ------------------------------------------------------------------------------
\section{The chi-squared distribution}
\label{sec:chi2}

\begin{definition}
If $s_1, \dots, s_k$ are independent standard normal variables, $\gamma = \sum_{i=1}^k s_i^2$ has the
$\chi^2$ distribution with $k$ degrees of freedom, written $\gamma \sim \chi^2_k$. It has mean $k$
and variance $2k$.
\end{definition}

\begin{theorem}[Mahalanobis distance of a Gaussian]
\label{thm:mahalanobis-chi2}
If $x \sim \N(\mu, \Sigma)$ in $\R^k$ with $\Sigma \succ 0$, then $(x-\mu)\T\Sigma\inv(x - \mu) \sim \chi^2_k$.
\end{theorem}
\begin{proof}
With $\Sigma = LL\T$, the whitened vector $s = L\inv(x - \mu)$ is $\N(0, I)$ by \cref{prop:affine}, and
$(x - \mu)\T\Sigma\inv(x-\mu) = s\T s = \sum s_i^2$.
\end{proof}

\begin{figure}[t]
  \centering
  \includegraphics[width=0.84\textwidth]{fig_chi2.pdf}
  \caption{$\chi^2$ densities and their 95\% quantiles. A 2D image measurement whose normalized
  innovation squared exceeds $5.99$ would occur by chance only 5\% of the time if the filter's
  covariance were correct.}
  \label{fig:chi2}
\end{figure}

This theorem is the basis of two tools. \emph{Gating} (\cref{ch:updates}) rejects a measurement
whose normalized innovation squared $r\T S\inv r$ exceeds the 95\% (or 99\%) quantile of
$\chi^2_{\dim r}$. \emph{Consistency testing} (\cref{ch:consistency}) checks that the normalized
estimation error squared behaves like a $\chi^2$ variable over time and over Monte Carlo runs.
Because independent $\chi^2$ variables add (degrees of freedom add), the sum of $N$ independent
NEES values of an $n$-dimensional state is $\chi^2_{Nn}$. That gives exact confidence bounds for
averages. \Cref{app:tables} tabulates the quantiles.

% ------------------------------------------------------------------------------
\section{Continuous-time noise}
\label{sec:ct-noise}

IMU noise is specified in continuous time, and getting its units right is essential.

\begin{definition}[Continuous-time white noise]
A zero-mean process $w(t) \in \R^n$ is white with \emph{power spectral density} (PSD) $Q_c$ if
$\E[w(t)\,w(s)\T] = Q_c\,\delta(t - s)$. For a scalar process we write $Q_c = \sigma^2$ and call
$\sigma$ the \emph{noise density}.
\end{definition}

White noise is an idealization with infinite variance. Its meaningful properties are those of its
integrals. If $w$ is measured in some unit $[u]$, then $\sigma$ has units $[u]/\sqrt{\mathrm{Hz}}
= [u]\sqrt{\mathrm{s}}$, because $\delta(t-s)$ has units of $1/\mathrm{s}$. For a gyroscope,
$[u] = \mathrm{rad/s}$, so $\sigma_g$ is in $\mathrm{rad/s/\sqrt{Hz}}$, exactly the units of
\code{gyroscope\_noise\_density} in EuRoC's \code{sensor.yaml}.

\begin{proposition}[Integrated white noise]
\label{prop:integrated-noise}
Let $w$ be scalar white noise with density $\sigma$. The $m$-fold integral
$I_m(t) = \int_0^t \int_0^{t_1}\cdots\int_0^{t_{m-1}} w(t_m)\,\dd t_m \cdots \dd t_1$ has zero mean and
\begin{equation}
  \Var\big(I_m(t)\big) = \sigma^2\,\frac{t^{2m-1}}{(2m-1)\,\big((m-1)!\big)^2}.
  \label{eq:integrated-noise}
\end{equation}
For $m = 1,2,3,4$ this gives $\sigma^2 t$, $\sigma^2 t^3/3$, $\sigma^2 t^5/20$ and $\sigma^2 t^7/252$.
\end{proposition}
\begin{proof}
By Cauchy's formula for repeated integration, $I_m(t) = \int_0^t \frac{(t - s)^{m-1}}{(m-1)!}\,w(s)\,\dd s$.
For a linear functional $\int g(s)w(s)\,\dd s$ of white noise, the variance is
$\iint g(s)g(s')\,\sigma^2\delta(s - s')\,\dd s\,\dd s' = \sigma^2\int g(s)^2\,\dd s$. Here
$\int_0^t (t-s)^{2m-2}\,\dd s/((m-1)!)^2 = t^{2m-1}/((2m-1)((m-1)!)^2)$.
\end{proof}

The single integral $I_1$ is the \emph{Wiener process} or \emph{random walk}, whose standard deviation
grows as $\sqrt t$. A bias modeled as $\dot b = w$ is a random walk. Integrating accelerometer white
noise twice into position gives $\sigma_a^2 t^3/3$, the analytic check of Milestone 5. Each further
integration adds two powers of $t$ to the variance. \Cref{ch:errordyn} uses this proposition to
produce the drift laws of \cref{fig:cov-growth}.

\paragraph{From continuous to discrete noise.} A digital sensor reports samples every $\Dt$. If each
sample is the average of the white-noise input over its interval (which an anti-aliasing filter
approximates), the sample noise has variance
\begin{equation}
  \Var\Big(\frac1\Dt\int_0^\Dt w(s)\,\dd s\Big) = \frac{\sigma^2}{\Dt},
  \qquad \sigma_{\mathrm{discrete}} = \frac{\sigma}{\sqrt{\Dt}}.
  \label{eq:discrete-noise}
\end{equation}
For EuRoC's gyroscope, $\sigma_g = 1.6968\times10^{-4}~\mathrm{rad/s/\sqrt{Hz}}$ at 200~Hz gives a
per-sample standard deviation of $2.4\times10^{-3}$~rad/s. Conversely, a random walk
$\dot b = w$ discretized over $\Dt$ gains variance $\sigma_b^2\Dt$ per step:
$b_{k+1} = b_k + \sqrt{\Dt}\,\sigma_b\, s_k$, $s_k \sim \N(0,1)$. The two conversions go in opposite
directions; mixing them up is a classic tuning error.

\section{Consistency}

\begin{definition}[Consistent estimator]
An estimator that reports $(\hat x, P)$ is \emph{consistent} if its error $e = x - \hat x$ satisfies
$\E[e] = 0$ and $\E[ee\T] = P$. It is \emph{overconfident} if the actual error covariance exceeds
$P$, and \emph{underconfident} if it is smaller.
\end{definition}

Consistency matters more than it first appears. An overconfident filter trusts its own prediction
too much, so it gives new measurements too little weight. Its errors then grow, which makes it
even more inconsistent, and it eventually rejects valid measurements in its gate. Most catastrophic
filter failures in practice begin as gradual overconfidence. \Cref{ch:consistency} develops the
tests and \cref{ch:observability} explains the structural cause in VIO.

\begin{takeaways}
\begin{itemize}
  \item Affine maps propagate covariances as $A\Sigma A\T$. This single fact drives prediction,
        cloning and stacking.
  \item In covariance form, marginalization deletes rows and columns and conditioning takes a Schur
        complement. Information form is the reverse.
  \item For linear-Gaussian models, the Bayesian posterior mean is the weighted least-squares
        solution.
  \item Linearization approximates a nonlinear transformation by a Gaussian, and fails when
        uncertainty is large relative to curvature.
  \item The Mahalanobis distance of a Gaussian vector is $\chi^2$ with $\dim$ degrees of freedom.
  \item White-noise densities have units of $[u]/\sqrt{\mathrm{Hz}}$. The $m$-fold integral of white
        noise has variance growing like $t^{2m-1}$.
\end{itemize}
\end{takeaways}

\begin{exercises}
  \item Let $x \sim \N(0, \diag(1, 4))$ and $y = Ax$ with $A$ a rotation by $30^\circ$. Compute
        $\Cov(y)$ and sketch the 1$\sigma$ ellipses of $x$ and $y$.
  \item Prove that $\Sigma_{a|b} \preceq \Sigma_{aa}$: conditioning never increases uncertainty in the
        Loewner order. When is the inequality an equality?
  \item Derive the information-form marginal $\Lambda_{aa} - \Lambda_{ab}\Lambda_{bb}\inv\Lambda_{ba}$
        from \cref{prop:block-inverse}, and show it equals $\Sigma_{aa}\inv$.
  \item Solve \eqref{eq:map-ls} by setting the gradient to zero, and show that the solution can be
        written $\hat x = \mu + K(z - H\mu)$ for $K = (P\inv + H\T R\inv H)\inv H\T R\inv$.
  \item For $x \sim \N(0, \sigma^2)$ and $y = x^2$, compute the exact mean of $y$ and its first-order
        linearized prediction. Explain the discrepancy.
  \item Verify \cref{prop:integrated-noise} by Monte Carlo for $m = 2$ and $m = 3$, simulating at
        200~Hz for 10~s with 1000 runs.
  \item An accelerometer has density $\sigma_a = 2\times10^{-3}~\mathrm{m/s^2/\sqrt{Hz}}$. A user
        samples it at 1~kHz instead of 200~Hz. How does the per-sample noise change? How does the
        position uncertainty after 10~s of integration change? Explain why.
  \item Show that the 95\% quantile of $\chi^2_k$ divided by $k$ tends to $1$ as $k\to\infty$, and
        explain what that implies for gating high-dimensional stacked measurements.
\end{exercises}
